\subsection{克利福德群。}

在本小节中，我们假设 A 是一个体，E 具有有限维数 m，且 Q 非退化。我们把 E 与它在 C(Q) 中的典范像等同起来。

定义 2. --- 我们把 C(Q)（相应地，C \(^{+}\) (Q)）中满足 sEs \(^{-1}\) = E 的可逆元素 s 所成的乘法群称为 Q 的克利福德群（相应地，Q 的特殊克利福德群）。

在本小节中，我们把 Q 的克利福德群和特殊克利福德群分别记作 G 与 G⁺。显然有 G⁺ = G ∩ C⁺(Q)。

定理 4. --- 对 \(s \in G\) 和 \(x \in E\) ，令 \(\varphi(s). x = sxs^{-1}\) 。

\begin{enumerate}
\def\labelenumi{\alph{enumi})}
\item
  映射 \(\varphi\) 是从 G 到 Q 的正交群 \(\mathbf{O}(\mathrm{Q})\) 中的同态，其核是 C(Q) 的中心 Z 中可逆元素所成的集。
\item
  集 \(E \cap G\) 是 E 的非奇异向量所成的集；当 \(x \in E \cap G, -\varphi(x)\) 就是关于与 x 正交的超平面的对称（变换）。
\item
  若 dim(E) 为偶数，则有 \(\varphi(G) = \mathbf{O}(Q)\) ，且 \(\varphi(G^{+})\) 在 \(\mathbf{O}(Q)\) 中的指数为 2（当 \(E \neq \{0\}\) 时），并等于 \(\mathbf{SO}(Q)\) （当 A 的特征为 \(\neq 2\) 时）。
\item
  若 dim(E) 为奇数（这蕴含 A 的特征 ≠ 2），则有 φ(G) = φ(G⁺) = SO(Q)。
\end{enumerate}

实际上，我们有 \(\mathrm{Q}(sxs^{-1}) = (sxs^{-1})^2 = sx^2s^{-1} = \mathrm{Q}(x)\) （对 \(s \in G\) 和 \(x \in E\)），这表明 \(\varphi(s) \in \mathbf{O}(Q)\) 。要使 \(\varphi(s) = 1\) ，必须且只需 s 与 E 的元素交换，亦即属于 C(Q) 的中心 Z。这就证明了 a)。

要使 E 的元素 \(x\) 属于 G，它必须可逆，也就是说，必须是非奇异向量（因为 \(x^2 = Q(x)\) ）。若确实如此，我们有 \(x^{-1} = Q(x)^{-1}x\) ，从而对每个 \(y \in E\) ，有

\[
x y x ^ {- 1} = \mathrm{Q} (x) ^ {- 1} x y x = \mathrm{Q} (x) ^ {- 1} x (\Phi (x, y) - x y) = - (y - \Phi (x, y) \mathrm{Q} (x) ^ {- 1} x)
\]

这就证明了 b)（§ 6，第 4 小节）。

引理 5. --- G 的每个元素 s 都具有形式 \(zs'\) ，其中 z 是 Z 的一个可逆元素，而 \(s'\) 属于 G ∩ C \(^{+}\) (Q) 或 G ∩ C \(^{-}\) (Q)；子群 G \(^{+}\) 在 G 中的指数为 2（当 E ≠ \{0\} 时）。

第二个断言显然由第一个推出，因为非奇异向量属于 \(\mathrm{G}\cap\mathrm{C}^{-}(\mathrm{Q})\) 。先设 dim(E) 为偶数，令 \(s=t'+t''\) ，其中 \(t'\in\mathrm{C}^{+}(\mathrm{Q})\)，\(t''\in\mathrm{C}^{-}(\mathrm{Q})\)；按定义有 \(sx=(\varphi(s).x)s\) （对一切 \(x\in E\)）；由于 \(t'x\) 与 \((\varphi(s).x)t'\)（相应地，\(t''x\) 与 \((\varphi(s).x)t'')\)）是奇（相应地，偶）元素，故有 \(t'x=(\varphi(s).x)t'\) ，从而有 \(s^{-1}t'x=xs^{-1}t'\) （对一切 \(x\in E\)）。由此得出 \(s^{-1}t'\in Z\) ，又因 dim(E) 为偶数，有 \(Z=A\) （第 4 小节定理 2 的推论），于是 \(t'=as\) ，其中 \(a\in A\) 。若 \(a\neq0\) ，则有 \(s=a^{-1}t'\) 且 \(t'\in G\cap C^{+}(Q)\) ；若 a=0，则 \(s=t''\in G\cap C^{-}(Q)\)，引理在这种情形得证。若 dim(E) 为奇数，则 A 的特征为 \(\neq2\) ，于是对每个 \(s\in G\) ，\(\varphi(s)\) 都是关于非奇异向量 \(x_{i} (i = 1,2,\ldots,h)\) 的诸对称（变换）之积（ \(\S 6, n^{o} 4\) ，命题 5）；若令 \(s' = x_{1}x_{2}\ldots x_{h}\) ，则有 \(\varphi(s) = \varphi(s')\) ，从而 \(s = zs'\) ，其中 \(z \in \mathbf{Z}\) ，且 \(s'\) 属于 \(C^{+}(Q)\) 或 \(C^{-}(Q)\)，视 \(h\) 为偶数或奇数而定。

设 dim(E) 为偶数。由于每个元素 \(u\) （\(\mathbf{O}(Q)\) 中的）都以唯一的方式扩张为 C(Q) 的自同构 \(\overline{u}\) （命题 1），且 C(Q) 是中心单的（定理 2），故 \(\overline{u}\) 是内自同构（第八章 § 10 第 1 小节定理 1）。因此存在 G 的元素 \(s\) 使 \(\varphi(s) = u\) 。另一方面，C(Q) 的中心包含于 C \(^{+}\) ，这蕴含 \(\varphi(G)/\varphi(G^{+})\) 同构于 G/G \(^{+}\) ，从而 \(\varphi(G^{+})\) 在 \(\varphi(G) = \mathbf{O}(Q)\) 中的指数为 2（当 E ≠ \{0\} 时）。这就证明了 \(c\) ) 的前两个断言。

最后，设 A 的特征 ≠ 2。这时 O(Q) 的每个元素 u 都是关于与非奇异向量 \(x_{i} (i = 1, \ldots, h) (\S 6, n^{o} 4, \text{prop. } 5)\) 正交的超平面的诸对称（变换）之积；因此有 \(u = (-1)^{h} \varphi(x_{1} \ldots x_{h})\) 且 det \((u) = (-1)^{h}\) 。要使 u 属于 SO(Q)，必须且只需 h 为偶数，这表明 \(\varphi(G^{+}) \supset \mathbf{SO}(Q)\) 。由于当 \(E \neq \{0\}\) 时 SO(Q) 在 O(Q) 中的指数为 2，故若 E 为偶数维，则有 \(\varphi(G^{+}) = \mathbf{SO}(Q)\) ，这便完成了 c) 的证明。另一方面，若 E 的维数为奇数，则 \(\varphi(G)\) 不包含正交变换 \(x \to -x\) ；事实上，后者扩张为 C(Q) 的主自同构 \(\alpha\) （第 1 小节），而由于 C(Q) 的中心 Z 含有 C⁻(Q) 的一个非零元素（定理 3），\(\alpha\) 不是内自同构。于是有 \(\varphi(G) \neq \mathbf{O}(Q)\) ，又由于 \(\varphi(G) \supset \varphi(G^{+}) \supset \mathbf{SO}(Q)\) 且 SO(Q) 在 O(Q) 中的指数为 2，故有 \(\varphi(G) = \varphi(G^{+}) = \mathbf{SO}(Q)\) 。这就证明了 d)。证毕。

子群 \(\varphi(G^{+})\) （\(\mathbf{O}(Q)\) 中的，当 \(E \neq \{0\}\) 时其指数为 2）称为 E 的旋转群，其元素称为旋转；它记作 \(\mathbf{O}^{+}(Q)\) 。注意，若 A 的特征不是 2，则有 \(\mathbf{O}^{+}(Q) = \mathbf{S}\mathbf{O}(Q)\) （参见习题 9）。

命题 4. — 设 β 是 C(Q) 的主反自同构（第 1 小节）。对每个 \(s \in G^{+}\) ，\(\beta(s)s\) 是一个标量。映射 N ： \(s \to \beta(s)s\) 是从 \(G^{+}\) 到 A 中非零元素所成的乘法群 A* 中的同态。

实际上，对 \(s \in G^{+}\) ，有 \(sEs^{-1} = E\) ，从而 \(\beta(s)^{-1}E\beta(s) = E\) ，这表明 \(\beta(s) \in G^{+}\) 。由于有 \(sx = (\varphi(s).x)s\) （对一切 \(x \in E\)），故有 \(x\beta(s) = \beta(sx) = \beta(s)(\varphi(s).x)\) ，因而 \(\beta(s)sx = \beta(s)(\varphi(s).x)s = x\beta(s)s\) ，这蕴含 \(\beta(s)s\) 属于 C(Q) 的中心。又由于 \(\beta(s)s\) 属于 C+(Q)，故 \(\beta(s)s\) 是一个标量（定理 2 与定理 3）。最后有 \(\beta(st)st = \beta(t)\beta(s)st = \beta(s)s\beta(t)t\) ，即对 G+ 中的 s、t，有 N(st) = N(s)N(t)。证毕。

标量 \(N(s) = \beta(s)s (s \in G^{+})\) 称为 s 的旋量范数。我们把同态 N 的核记作 \(G_{0}^{+}\) ，并称之为约化克利福德群。像 \(\varphi(G_{0}^{+})\) 记作 \(\mathbf{O}_{0}^{+}(Q)\) ，称为 Q 的约化正交群。由于 \(\varphi\) 在 \(G^{+}\) 上的限制的核是 C(Q) 的中心中偶可逆元素所成的集（定理 4），从而可与 A* 等同（定理 2 与定理 3），故 \(\varphi(G^{+})/\mathbf{O}_{0}^{+}(Q)\) 同构于 \(G^{+}/A^{*}G_{0}^{+}\) ，从而也同构于 \(N(G^{+})/N(A^{*})\) ，特别地它是交换的。显然 \(N(A^{*})\) 是 A* 中元素的平方所成的子群 \((A^{*})^{2}\) 。若 Q 的指数 \textgreater0，则对每个 \(a \in A^{*}\) ，存在 E 的两个元素 x 与 y，使得 \(Q(x) = a\) 且 \(Q(y) = 1 (\S 4, n^{\circ} 2, \text{prop. } 4)\) ；由于 \(xy \in G^{+}\) 且 \(N(xy) = Q(x)Q(y) = a\) ，这表明 \(N(G^{+}) = A^{*}\) ，从而 \(\varphi(G^{+})/\mathbf{O}_{0}^{+}(Q)\) 同构于 \(A^{*}/(A^{*})^{2}\) 。

习题. --- ¶ 1) 当 E₁ 与 E₂ 是 E 的任意两个互补子模时，证明第 3 小节定理 1 的推论 3 与推论 4。（先证推论 4：为此，证明张量积 C(Q₁) ⊗ C(Q₂) 依推论 4 的陈述中所作的符号约定具有一个代数结构，且这个代数 S 与 C(Q) 解决同一个泛映射问题；为此，对从 E 到 A 上代数 D 中的、满足 (f(x))² = Q(x)·1 的每个线性映射 f，考虑 C(Qᵢ) 到 D 中的同态 \(\bar{f}_i\) ，它满足 \(\bar{f}_i(1) = 1\) ，且 \(\bar{f}_i(\rho_{Q_i}(x_i)) = f(x_i)\) 对 \(x_i \in \mathrm{E}_i\) ( \(i = 1, 2\) ) 成立，并证明存在代数 S 到 D 中的同态 \(\bar{f}\) 使得

\[
\bar {f} (z _ {1} \otimes z _ {2}) = \bar {f} _ {1} (z _ {1}) \bar {f} _ {2} (z _ {2})
\]

对 \(z_{i}\in\mathrm{C}(\mathrm{Q}_{i})\) ( \(i=1,2\) )。然后为证推论 3，考虑作为 Q \(_{1}\) 与 Q \(_{2}\) 的外直和的二次型 Q'（§ 3，第 4 小节），并注意有 Q'(x)=Q(x)+F(x,x)，其中 F 是由 F(x \(_{1}\) +x \(_{2}\) , y \(_{1}\) +y \(_{2}\) ) = -Φ(x \(_{1}\) , y \(_{2}\) )（对 x \(_{i}\) 、y \(_{i}\) 属于 E \(_{i}\) ( \(i=1,2\) )）定义的双线性型；然后应用命题 3。）

¶ 2) 设 E 是两个正交子模 E \(_{1}\) 、E \(_{2}\) 的直和，并用 Q \(_{i}\) 表示 Q 在 E \(_{i}\) 上的限制（i = 1, 2）；

再设存在 \(u \in \mathrm{C}^{+}(\mathrm{Q}_{2})\) 满足 \(u^{2} = a.1\) （其中 \(a\) 在 A 中可逆），且 \(u\rho_{\mathrm{Q}_{2}}(x_{2}) = -\rho_{\mathrm{Q}_{2}}(x_{2})\) 对一切 \(x_{2} \in \mathrm{E}_{2}\) 成立。证明存在一个同构 \(\varphi\) ，它把 C(Q) 映到张量积 \(\mathrm{C}(a\mathrm{Q}_{1}) \otimes \mathrm{C}(\mathrm{Q}_{2})\) 上（张量积按第三章 § 3 第 1 小节所定义），且具有如下性质：对每个 \(x = x_{1} + x_{2}\) ( \(x_{i} \in \mathrm{E}_{i}, i = 1, 2\) )，有

\[
\varphi (\rho_ {q} (x)) = \rho_ {a q _ {1}} (x _ {1}) \otimes u ^ {- 1} + 1 \otimes \rho_ {q _ {2}} (x _ {2}).
\]

（由 C(Q) 的泛性质推出同态 \(\varphi\) 的存在性。另一方面，有一个同态 \(g_{1}\) （从 C(aQ \(_{1}\) ) 到 C(Q) 中），满足 \(g_{1}(\rho_{aQ_{1}}(x_{1})) = h_{2}(u)\rho_{Q_{1}}(x_{1})\) ，其中 \(h_{2}\) 表示 C(Q \(_{2}\) ) 到 C(Q) 中的典范同态；然后注意 \(g_{1}(z_{1})\) 与 \(h_{2}(z_{2})\) 对 \(z_{1} \in\) C(aQ \(_{1}\) ) 与 \(z_{2} \in\) C(Q \(_{2}\) ) 交换，由此推出 \(\varphi\) 的逆同态的存在性。）

当 A 是特征 ≠ 2 的体时，这一结果适用于哪些情形（利用定理 2 的推论之后的注记 2）？

\begin{enumerate}
\def\labelenumi{\arabic{enumi})}
\setcounter{enumi}{2}
\tightlist
\item
  \begin{enumerate}
  \def\labelenumii{\alph{enumii})}
  \tightlist
  \item
    在第 2 小节的记号下，设 \(x_{i}\) ( \(1 \leqslant i \leqslant n\) ) 是 E 中满足 \(f(x_{i}) = 0\) 的元素；证明有 \(i_{j}(x_{1} \otimes x_{2} \otimes \cdots \otimes x_{n}) = 0\) 。特别地，若 F 是 E 上的一个双线性型，满足 F( \(x_{i}, x_{j}\) ) = 0 （对 \(i > j\) ），则有 \(i_{x_{1}}^{\mathrm{F}}(x_{2} \otimes x_{3} \otimes \cdots \otimes x_{n}) = 0\) 。
  \end{enumerate}
\end{enumerate}

\begin{enumerate}
\def\labelenumi{\alph{enumi})}
\setcounter{enumi}{1}
\tightlist
\item
  在第 3 小节命题 3 的记号下，设 \(x_i\) ( \(1 \leqslant i \leqslant n\) ) 是 E 中满足 \(F(x_i, x_j) = 0\) （对 \(i > j\) ）的元素。证明我们有 \(\overline{\lambda}_{\mathrm{F}}(\rho_{\mathrm{Q}'}(x_1) \ldots \rho_{\mathrm{Q}'}(x_n)) = \rho_{\mathrm{Q}}(x_1) \ldots \rho_{\mathrm{Q}}(x_n)\) （利用 \(a\) ) 及第 2 小节的公式 (10)）。
\item
  假设 A 是特征 \(\neq 2\) 的体；对 E 上的每个二次型 Q，令 \(\mu_{\mathrm{Q}}\) 为映射 \(\overline{\lambda}_{\mathrm{F}}\) （从 C(Q) 到 \(\wedge\) E 上），它对应于 \(F(x, y) = \frac{1}{2}\Phi(x, y)\) 。证明：若诸向量 \(x_i\) ( \(1 \leqslant i \leqslant n\) ) 两两正交，则有
\end{enumerate}

\[
\mu_ {\mathrm{Q}} (x _ {1} x _ {2} \dots x _ {n}) = x _ {1} \wedge x _ {2} \wedge \dots \wedge x _ {n}.
\]

由此推出：若 \(y_{i}\) ( \(1 \leqslant i \leqslant n\) ) 是 \(n\) 个任意向量，则有

\[
n! \mu_ {0} ^ {- 1} (y _ {1} \wedge y _ {2} \wedge \dots \wedge y _ {n}) = \sum_ {\sigma \in \mathfrak {S} _ {n}} \varepsilon_ {\sigma} y _ {\sigma (1)} y _ {\sigma (2)} \dots y _ {\sigma (n)}.
\]

\begin{enumerate}
\def\labelenumi{\arabic{enumi})}
\setcounter{enumi}{3}
\tightlist
\item
  设 \(C_{h}\) 为 C(Q) 的子模，它由 \(\rho_{Q}(E)\) 的 k 个元素之积（\(0 \leqslant k \leqslant h\)）生成。证明映射
\end{enumerate}

\[
(x _ {1}, \dots , x _ {h}) \rightarrow \rho_ {\mathrm{Q}} (x _ {1}) \dots \rho_ {\mathrm{Q}} (x _ {h})
\]

在过渡到商之后定义了一个从 \(E^{h}\) 到 \(C_{h}/C_{h-1}\) 中的交错多重线性映射。由此得到一个线性映射 \(\pi_{h}\) （从 \(\bigwedge E\) 到 \(C_{h}/C_{h-1}\) 中）；证明当 E 是自由模时 \(\pi_{h}\) 是同构。若 f 是一个正交变换，我们有 \(C(f)\circ\pi_{h}=\pi_{h}\circ(\bigwedge f)\) 。当 A 是特征 \(\neq2\) 的体时，证明 \(\pi_{h}\) 是习题 3 c) 中所定义的映射 \(\mu_{Q}^{-1}\) 的限制。

\begin{enumerate}
\def\labelenumi{\arabic{enumi})}
\setcounter{enumi}{4}
\tightlist
\item
  在第 2 小节的记号下，考虑映射 \(i_{f}\) （从 \(\wedge\) E 到其自身）。证明有
\end{enumerate}

\[
i _ {f} \left(x _ {1} \wedge \dots \wedge x _ {h}\right) = \sum_ {i = 1} ^ {h} (- 1) ^ {i - 1} f \left(x _ {i}\right) \left(x _ {1} \wedge \dots \wedge x _ {i - 1} \wedge x _ {i + 1} \wedge \dots \wedge x _ {h}\right)
\]

并由此推出，\(i_{f}\) 不是别的，正是右内乘 \(z \to z \sqcup f\) （第三章 § 8 第 4 小节定义 2）。

\begin{enumerate}
\def\labelenumi{\arabic{enumi})}
\setcounter{enumi}{5}
\item
  证明：若 E 具有一个关于 Q 的正交基 \((e_{1}, e_{2})\) ，且令 \(\alpha_{i} = \mathrm{Q}(e_{i})\) （对 i = 1, 2），则代数 C(Q) 同构于 A 上与偶 \((\alpha_{1}, \alpha_{2})\) 相应的四元数代数。
\item
  设 A 为极大有序域，E 为 A 上维数为偶数 2r 的向量空间，Q 为 E 上正的或负的非退化二次型。证明：若 Q 是正的，则当 \(r(r-1)/2\) 为偶数时，C(Q) 同构于 A 上的一个矩阵代数；当 \(r(r-1)/2\) 为奇数时，同构于 A 上四元数代数上的一个矩阵代数。若 Q 是负的，则当 \(r(r+1)/2\) 为偶数时，C(Q) 同构于 A 上的一个矩阵代数；当 \(r(r+1)/2\) 为奇数时，同构于 A 上四元数代数上的一个矩阵代数（利用习题 2 与习题 6）。在这些不同情形下，C⁺(Q) 的结构如何？
\item
  设 A 为环 \(\mathbf{Z}/(4)\) ，E 为左 A-模 \(\mathbf{Z}/(2)\) ；若令 \(Q(0)=0\) ，且令 \(Q(u)=1\) （对 E 中唯一的元素 \(u\neq0\) ），则 Q 是 E 上的一个二次型。证明 A-模 C(Q) 与 \(\wedge\) E 不同构（可证明 C(Q) 同构于两个同构于 E 的模的直和）。
\end{enumerate}

¶ 9) 设 A 是特征 2 的体，E 是维数为有限偶数 2r 的向量空间，Q 是 E 上的非退化二次型。

\begin{enumerate}
\def\labelenumi{\alph{enumi})}
\tightlist
\item
  设 \((e_i)\) 是 E 关于与 Q 相伴的交错双线性型 \(\Phi\) 的一个辛基（§ 5，第 1 小节）。证明元素
\end{enumerate}

\[
z = e _ {1} e _ {2} + e _ {3} e _ {4} + \dots + e _ {2 r - 1} e _ {2 r}
\]

（C \(^{+}\) (Q) 中的）与单位元一起构成 C \(^{+}\) 的中心 Z 的一个基。要使 Z 是两个域的直和，必须且只需元素

\[
\Delta (Q) = Q \left(e _ {1}\right) Q \left(e _ {2}\right) + Q \left(e _ {3}\right) Q \left(e _ {4}\right) + \dots + Q \left(e _ {2 r - 1}\right) Q \left(e _ {2 r}\right)
\]

（称为 Q 关于辛基 \((e_{i})\) 的伪判别式）具有形式 \(\lambda^{2} + \lambda\) ，其中 \(\lambda \in A\) 。

\begin{enumerate}
\def\labelenumi{\alph{enumi})}
\setcounter{enumi}{1}
\tightlist
\item
  设 \(u\) 是关于 \(\Phi\) 的一个相似变换（ \(\S 6\) ，第 5 小节），其乘子为 \(\mu(u)\) ，并令 \(Q_{1}(x) = Q(u(x))\) 。我们置
\end{enumerate}

\[
u (e _ {2 i - 1}) = \sum_ {j = 1} ^ {r} a _ {i j} e _ {2 j - 1} + \sum_ {j = 1} ^ {r} b _ {i j} e _ {2 j},
\]

\[
u (e _ {2 i}) = \sum_ {j = 1} ^ {r} c _ {i j} e _ {2 j - 1} + \sum_ {j = 1} ^ {r} d _ {i j} e _ {2 j},
\]

以及 \(\mathrm{Q}(e_{2i - 1}) = \alpha_i, \mathrm{Q}(e_{2i}) = \beta_i (1 \leqslant i \leqslant r)\) 。证明我们有

\[
\Delta (\mathrm{Q} _ {1}) = (\mu (u)) ^ {2} \Delta (\mathrm{Q}) + (\mathrm{D} (u)) ^ {2} + \mu (u) \mathrm{D} (u)
\]

其中

\[
\mathbf {D} (u) = \sum_ {i, j} (\alpha_ {j} a _ {i j} c _ {i j} + \beta_ {j} b _ {i j} d _ {i j} + b _ {i j} c _ {i j})
\]

（u 关于基 \((e_{i})\) 的迪克森不变量）。（注意元素

\[
(\mu (u)) ^ {- 1} (u (e _ {1}) u (e _ {2}) + u (e _ {3}) u (e _ {4}) + \dots + u (e _ {2 r - 1}) u (e _ {2 r}))
\]

属于 Z。）要使 \(u\) 是 Q 的一个相似变换（§ 4，习题 9）且其乘子为 \(\mu(u)\) ，必须且只需 D(u) = 0 或 D(u) = \(\mu(u)\) ；Q 的满足 D(u) = 0 的相似变换称为正相似变换。

\begin{enumerate}
\def\labelenumi{\alph{enumi})}
\setcounter{enumi}{2}
\tightlist
\item
  证明：若 v 是关于 \(\Phi\) 的一个相似变换，u 是 Q 的一个相似变换，则有
\end{enumerate}

\[
\mathrm{D} (u v) = \mu (v) \mathrm{D} (u) + \mu (u) \mathrm{D} (v)
\]

（考虑 \(\varrho\) 关于由诸 \(u(e_{2i-1})\) 与 \((\mu(u))^{-1}u(e_{2i})\) （对 \(1 \leqslant i \leqslant r\)）所构成的辛基的迪克森不变量）。由此推出，Q 的正相似变换构成 Q 的相似变换群中指数为 2 的一个正规子群。

\begin{enumerate}
\def\labelenumi{\alph{enumi})}
\setcounter{enumi}{3}
\item
  若 \(u\) 是关于与 E 中一个非奇异向量正交的超平面的对称（变换），证明 \(\mathrm{D}(u)=1\) 。由此推出，群 \(\varphi(\mathrm{G}^{+})\) （第 5 小节的记号）就是 § 6 习题 28 c) 中定义的群 \(\mathbf{SO}(\mathrm{Q})\) 。
\item
  设 \(u \in \mathbf{O}(Q)\) ，并设 \(\Lambda\) 代数闭。证明：要使 \(u \in \mathbf{SO}(Q)\) ，必须且只需模 \(E_u\) 的初等因子的个数为偶数。（在 § 5 习题 15 的记号下，首先注意，若 \(p, q\) 是 \(u\) 的极小多项式的两个不同的不可约因子，则 \(G(p, q)\) 所分解成的诸不可分解子模的个数等于 \(G(q, p)\) 所分解成的诸不可分解子模的个数。另一方面注意，\(G(p, p) = \{0\}\) 除去 \(p = X - 1\) 的情形。最后证明可化归到 \(E_u\) 等于 \(G(p, p)\) （其中 \(p = X - 1\) ）且不可分解的情形。这时存在 E 的一个辛基 ( \(e_i\) )，使得 \(e_1, e_3, \ldots, e_{2k-1}\) 构成 \((p(u))^{2r-k}(E)\) 的一个基（对 \(1 \leqslant k \leqslant r\) ）；证明 \(e_1, e_3, \ldots, e_{2r-3}\) 都是奇异向量，并由此断定 \(D(u) = 1\) 。）
\end{enumerate}

¶ 10) 设 A 是一个体，E 是有限维向量空间，Q 是 E 上的退化二次型；设 M 是 E 中 E \(^{0}\) 的一个补子空间，B 是 Q 在 M 上的限制的（半单）克利福德代数。

\begin{enumerate}
\def\labelenumi{\alph{enumi})}
\item
  先设 A 的特征 ≠ 2。设 L 是 Q 在 E⁰ 上的限制的克利福德代数（同构于 ∧ E⁰），R₀ 是它的根基（由 E⁰ 在 L 中生成的理想，在 L 中余维数为 1）；证明 C(Q) 的根基 R （在相差一个同构的意义下）是这样得到的：按定理 1 的推论 4 中的方式在 \(B \otimes_{A} R_{0}\) 上定义代数结构；并证明 C(Q)/R 同构于 B，且 C(Q) 是 B 与 R 的直和。
\item
  设 A 的特征为 2。设 F 是 \(\mathbf{E}^0\) 中由奇异向量 \(x \in \mathbf{E}^0\) 组成的子空间，N 是 F 在 \(\mathbf{E}^0\) 中的一个补。若 \((a_i)_{1 \leqslant i \leqslant d}\) 是 N 的一个基，且 \(\mathrm{Q}(a_i) = \alpha_i\) ，则诸元素 \(\alpha_i^{1/2}\) （取在 A 的一个代数闭包中）在 A 上线性无关。设 \((\alpha_i^{1/2})_{1 \leqslant i \leqslant e}\) 是域 \(\mathrm{A}_1 = \mathrm{A}(\alpha_1^{1/2}, \ldots, \alpha_d^{1/2})\) 在 A 上的一个 2-基（第五章 § 8，习题 1），并令 \(h = \dim F\) 。若 \(\mathrm{B}_1\) 是中心单代数 \(\mathrm{B} \otimes_{A} \mathrm{A}_1\) ，证明 \(\mathrm{C(Q)}\) 同构于代数 \(\mathrm{B}_1 \otimes_{A_1} \mathrm{L}_1\) ，其中 \(\mathrm{L}_1\) 是一个维数为 \(h + d - e\) 的向量空间在 \(\mathrm{A}_1\) 上的外代数。若 \(\Re_1\) 是 \(\mathrm{L}_1\) 的根基（余维数为 1（在 \(\mathrm{A}_1\) 上），于 \(\mathrm{L}_1\) 中），则根基 \(\Re\) （\(\mathrm{C(Q)}\) 的）同构于代数 \(\mathrm{B}_1 \otimes_{A_1} \Re_1\) ，\(\mathrm{C(Q)/R}\) 同构于 \(\mathrm{B}_1\) ，且 \(\mathrm{C(Q)}\) 是 \(\mathrm{B}_1\) 与 \(\Re\) 的一个直和。
\item
  假设仍如 \(b\) ；再设 \(h=0\) ， \(d=1\) ， \(e=0\) （这蕴含 dim M 为偶数）。若 Q \(_{0}\) 是 Q 在 M 上的限制，证明 C \(^{+}\) (Q) 同构于 C(Q \(_{0}\) )。
\end{enumerate}

11) a) 在第 5 小节的假设与记号下，证明 N(G⁺) 是 A* 中由诸乘积 Q(x)Q(y) 生成的子群，其中 x 与 y 取遍 E 的非迷向向量所成的集。（化归到 A ≠ F₂ 的情形；利用命题 5 与 § 6 的习题 28 c)，以及 § 9 的习题 9 d)。）Q 的指数 ≥ 1 的情形。

\begin{enumerate}
\def\labelenumi{\alph{enumi})}
\setcounter{enumi}{1}
\item
  再假设 \(\mathbf{A} \neq \mathbf{F}_2\) ，E 的维数为 \(n \geqslant 2\) ，且 Q 的指数 \(>0\) 。证明 \(\mathbf{O}_0^+(\mathbf{Q})\) 是群 \(\mathbf{O}(\mathbf{Q})\) 的换位子群。（利用 § 6 的习题 17 c) 与 28 e) 化归到 \(n = 2\) 的情形。）
\item
  我们保持 \(b\) 的假设，再设 A 的特征 \(\neq 2\) 且 E 为偶数维。要使 E 的自同构 \(x \to -x\) 属于 \(\mathbf{O}_0^+(\mathbf{Q})\) ，必须且只需 Q 的判别式（关于 E 的任意一个基）是一个平方。
\end{enumerate}

设 a 是 A 的一个可逆元素；证明存在唯一的同构 \(\theta_{a}\) ，从 C \(^{+}\) (Q) 到 C \(^{+}\) (aQ) 上，使得

\[
\theta_ {a} (\rho (x) \rho (y)) = a ^ {- 1} \rho_ {1} (x) \rho_ {1} (y)
\]

对 E 中的一切 \(x, y\) （ \(\rho\) 与 \(\rho_{1}\) 分别表示 E 到 C(Q) 中与 C(aQ) 中的典范映射）。由此推出：若 A 是特征 \(\neq 2\) 的体，E 是奇数维向量空间，Q 是非退化二次型，则 C(Q) 与 C(aQ) 同构（参见习题 7）。

\begin{enumerate}
\def\labelenumi{\alph{enumi})}
\setcounter{enumi}{1}
\tightlist
\item
  设 A 是一个体，E 是维数为偶数 2r \textgreater{} 0 的向量空间，Q 是非退化二次型。设 u 是关于 Q 的一个相似变换；证明存在唯一的 A-自同构 \(\overline{u}\) （代数 \(\mathrm{C}^{+}(\mathrm{Q})\) 的），使得对 \(1 \leqslant h \leqslant r\) 与 \(x_{i} \in E (1 \leqslant i \leqslant 2h)\) 有
\end{enumerate}

\[
\overline {{{u}}} (x _ {1} x _ {2} \dots x _ {2 h}) = \mu^ {- h} u (x _ {1}) u (x _ {2}) \dots u (x _ {2 h})
\]

其中 \(\mu\) 表示 \(u\) 的乘子。要使 \(\bar{u}\) 是内自同构，必须且只需 \(u\) 是正相似变换（§ 6，第 5 小节及 § 9，习题 9 b)）。再设 \(r \geqslant 2\) ；这时，要使 \(\bar{u}\) 是恒同，必须且只需 \(u\) 是一个位似。

证明自同构 \(\bar{u}\) （C \(^{+}\) (Q) 的）是 C(Q) 的一个内自同构在 C \(^{+}\) (Q) 上的限制。

¶ 13) 设 A 是特征 ≠ 2 的体，E 是维数为偶数 2r \textgreater{} 0 的向量空间，Q 是非退化二次型。我们用 \(E_{h}\) 表示 \(\Lambda E\) 在习题 3 c) 的同构 \(\mu_{Q}\) 下的原像，因而 \(E_{h}\) 是 C(Q) 中由诸乘积 \(x_{1}x_{2}\ldots x_{h}\) 生成的子空间，其中诸 \(x_{i}\) ( \(1 \leqslant i \leqslant h\) ) 两两正交。

\begin{enumerate}
\def\labelenumi{\alph{enumi})}
\tightlist
\item
  设 \(\alpha\) 是 A 的一个元素，\(s\) 是 E \(_2\) 的一个元素。证明：若 \((\alpha + s)^2 \in \mathbf{A}\) ，则或者 \(s = 0\) ，或者 \(\alpha = 0\) 且 \(s = xy\) ，其中 \(x\) 与 \(y\) 是两个正交向量。（注意，若 \(t = \mu_Q(\alpha + s)\) ，则 \(t \wedge t\)
\end{enumerate}

必然属于 A + ∧ E，这可借助于 E 的一个正交基表达 s 看出；再借助 § 5 第 1 小节定理 1 的推论 2 推出，必然有 t = β + x ∧ y，其中 β ∈ A，x、y 在 E 中。

\begin{enumerate}
\def\labelenumi{\alph{enumi})}
\setcounter{enumi}{1}
\item
  设 \((x, y)\) 与 \((u, v)\) 是 E 中两对正交的非迷向向量，\(\mathrm{P}_{xy}\) 与 \(\mathrm{P}_{uv}\) 分别是平面 \(Ax + Ay\) 与 \(Au + Av\) 。要使在 C(Q) 中有 \((xy)(uv) = -(uv)(xy)\) ，必须且只需 \(\mathrm{P}_{xy} + \mathrm{P}_{uv}\) 是一个维数为 3 的非迷向子空间，且在其中 \(\mathrm{P}_{xy}\) 与 \(\mathrm{P}_{uv}\) 弱正交（§ 3，习题 11）。
\item
  设 g 是 \(\mathrm{C}^{+}(\mathrm{Q})\) 的一个 A-自同构，它把 \(E_{2}\) 映到其自身；证明存在关于 Q 的一个相似变换 u，使得 \(g = \bar{u}\) （习题 12 b)）。（设 \((e_{i})_{1 \leqslant i \leqslant 2r}\) 是 E 的一个正交基；利用 a) 与 b)，证明存在 E 的一个正交基 \((x_{i})\) ，使得 \(g(e_{1}e_{i}) = x_{1}x_{i}\) （对 \(2 \leqslant i \leqslant 2r\) ）。）
\end{enumerate}

\begin{enumerate}
\def\labelenumi{\arabic{enumi})}
\setcounter{enumi}{13}
\tightlist
\item
  设 A 是特征 \(\neq 2\) 的体，E 是 A 上维数为 \(n\) 的向量空间，Q 与 Q₁ 是 E 上两个非退化二次型；设 \(\tilde{\Delta}\) （相应地，\(\tilde{\Delta}_{1}\) ）是 Q 的判别式 \(\Delta\) （相应地，\(\Delta_{1}\) ，Q \(_{1}\) 的判别式）关于 E 的一个基在群 A \(^{*}\) /(A \(^{*}\) )² 中的类，这个类不依赖于 E 中所选的基。我们假设 \(n = 2\) 。
\end{enumerate}

\begin{enumerate}
\def\labelenumi{\alph{enumi})}
\item
  证明：若 \(\tilde{\Delta} = \tilde{\Delta}_{1}\) 且克利福德代数 C(Q) 与 C(Q \(_{1}\) ) 同构，则 Q 与 Q \(_{1}\) 等价（在 C(Q) 上考虑二次型 \(z \to z\bar{z}\) ，其中 \(z \to \bar{z}\) 是 C(Q) 的唯一的对合反自同构，其不变元所成的集是 C(Q) 的中心（参见习题 6 及第八章 § 11，习题 5 e)）；应用维特定理）。
\item
  要使 C(Q) 同构于矩阵代数 \(\mathbf{M}_{2}(\mathrm{A})\) ，必须且只需对 E 中至少一个 \(x \neq 0\) ，存在 E 中的 \(y\) 、\(z\) 使得 Q(x) + Q(y)Q(z) = 0；若确是这样，则这个性质对 E 中每个 \(x \neq 0\) 都成立。
\end{enumerate}

¶ 15) 我们保持习题 14 的记号与一般假设，但设 n = 3。

\begin{enumerate}
\def\labelenumi{\alph{enumi})}
\item
  证明 \(\mathrm{C}^{+}(\mathrm{Q})\) 同构于 A 上的一个四元数代数，且 \(\beta\) 是这个代数的反自同构 \(z\to \bar{z}\) ，其不变元所成的集是 \(\mathrm{C}^{+}(\mathrm{Q})\) 的中心；若 P 是 \(\mathrm{C}^{+}(\mathrm{Q})\) 中由纯四元数（即满足 \(z = -\bar{z}\) 的那些元素；第八章 § 11，习题 6）组成的子空间，则二次型 \(z\to z\bar{z}\) 在 P 上的限制等价于 \(\lambda Q\) ，其中 \(\lambda \in A\) 。由此推出，要使 \(C^{+}(Q)\) 是一个域，必须且只需 \(Q\) 的指数为 0。
\item
  证明：若 \(\tilde{\Delta} = \tilde{\Delta}_{1}\) 且克利福德代数 C(Q) 与 C(Q \(_{1}\) ) 同构，则 Q 与 Q \(_{1}\) 等价。（先考虑 - \(\Delta\) 是 K 中平方的情形，并证明这时 C \(^{+}\) (Q) 与 C \(^{+}\) (Q \(_{1}\) ) 同构；然后如习题 14 a) 那样论证，利用 a) 及第八章 § 11 的习题 6。一般情形利用习题 12 a)。）
\item
  证明（关于型 Q 的）特殊克利福德群 G \(^{+}\) 恒同于 C \(^{+}\) (Q) 的可逆元素所成的群。（若 ( \(e_{1}\) , \(e_{2}\) , \(e_{3}\) ) 是 E 的一个正交基，且在 C(Q) 中 \(j = e_{1}e_{2}e_{3}\) ，则注意 \(x \to xj\) 是从 E 到 P 上的一个向量空间同构。）
\item
  由 a) 与 c) 推出：若 Q 的指数为 1，则旋转群 \(\mathbf{O}^{+}(\mathrm{Q})\) 同构于射影群 \(\mathbf{PGL}_{2}(\mathrm{A})\) （第二章 \(2^{\mathrm{e}}\) 版，附录 III，第 6 小节）。
\end{enumerate}

我们保持习题 14 的一般假设与记号，但设 n = 4。

\begin{enumerate}
\def\labelenumi{\alph{enumi})}
\item
  给出一个例子，其中 \(\tilde{\Delta} = \tilde{\Delta}_{1}\) ，\(C(Q)\) 与 \(C(Q_{1})\) 同构，但 \(Q\) 与 \(Q_{1}\) 不等价（参见习题 7）。
\item
  设 \((e_i)_{1\leqslant i\leqslant 4}\) 是 E 关于 Q 的一个正交基，\(Q_0\) 是 Q 在超平面 H = Ae \(_1\) + Ae \(_2\) + Ae \(_3\) 上的限制。证明：若 Z 是 C \(^+\) (Q) 的中心，则代数 C \(^+\) (Q) 同构于张量积 Z ⊗ \(_A\) C \(^+\) (Q \(_0\) )。对每个 z ∈ C \(^+\) (Q)，有 β(z)z ∈ Z ；要使 z 属于特殊克利福德群 G \(^+\) ，必须且只需 z 可逆且 β(z)z ∈ A.1。由此推出，群 O \(_0^+\) (Q) 同构于下述群对 \{1,-1\} 的商群，该群由满足 β(z)z = 1 的元素 z ∈ Z ⊗ \(_A\) C \(^+\) (Q \(_0\) ) 组成。
\item
  我们假设 \(\Delta\) 不是 A 中的平方（由维特定理，这蕴含 Q 的指数为 0 或 1）。若 \(Q_0'\) 是由 \(Q_0\) 经标量域扩张到 \(A' = A(\sqrt{\Delta})\) 而得的二次型，则由 \(b\) ) 推出 \(\mathbf{O}_0^+(Q)\) 同构于 \(\mathbf{O}_0^+(Q_0')\) 。特别地，若 Q 的指数为 1，则 \(\mathbf{O}_0^+(Q)\) 是 \(\mathbf{O}(Q)\) 的换位子群，且同构于 \(\mathbf{PSL}_2(A')\) （参见习题 15 \(d\) ) 及第三章 § 7，习题 8）。
\item
  我们假设 \(\Delta\) 是 A 中的平方（由维特定理，这蕴含 Q 的指数为 0 或 2），且 \(Q(e_4) = 1\) 。若令 \(j = e_1e_2e_3\) ，则 \(x \in E\) 可用唯一的方式写成 \(x = \alpha e_4 + jz\) ，其中 \(\alpha \in A\) ，而 \(z\) 是 \(L = C^+(Q_0)\) 中的一个纯四元数（习题 15 a)）；若令 \(\psi(x) = \alpha + z\) ，则 \(\psi\) 是从 E 到 L 上的一个向量空间同构。设 \(Z = Ac' + Ac''\) ，其中 \(c'\) 与 \(c''\) 是 Z 中两个正交的幂等元；每个可逆元素 \(s \in C^+(Q)\) 都可用唯一的方式写成 \(s = uc' + vc''\) ，其中 \(u\) 与 \(\nu\) 属于 L；要使 \(s\) 属于特殊克利福德群 \(G^+\) ，必须且只需 \(u\bar{u} = v\bar{v}\) ，而这时有 \(\psi(sxs^{-1}) = u\psi(x)v^{-1}\) （对一切 \(x \in E\) ）。由此推出，\(\mathbf{O}_0^+(Q)\) 对其中心（它是一个 2 元群，参见习题 11 c)）的商群同构于乘积 \(\mathbf{O}_0^+(Q_0) \times \mathbf{O}_0^+(Q_0)\) ；特别地，若 Q 的指数为 2，则这个商群同构于 \(\mathbf{PSL}_2(A) \times \mathbf{PSL}_2(A)\) 。
\end{enumerate}

\begin{enumerate}
\def\labelenumi{\arabic{enumi})}
\setcounter{enumi}{16}
\tightlist
\item
  设 K 是特征 ≠ 2 的域，A 是 K 上关于不定元的可数族的有理分式域 \(\mathrm{K}(\mathrm{X}_{n})_{n\in\mathbb{N}}\) （第四章 § 3，第 1 小节）。设 E 是 A 上的一个向量空间，具有可数基 \((e_{n})_{n\in\mathbb{N}}\) ，并设 Φ 是 E 上的一个对称双线性型，\((e_{n})\) 关于它是正交基，且 \(\Phi(e_{n}, e_{n}) = \mathrm{X}_{n}\) （对一切 \(n \in N\)）。若令 \(\mathrm{Q}(x) = \Phi(x, x)\) ，证明克利福德代数 C(Q) 是一个域（参见第八章 § 12，习题 14）。
\end{enumerate}

\section{角}

在本节中，A 表示一个特征 \(\neq 2\) 的域。设 E 是 A 上维数为 2 的向量空间，并设 \(\Phi\) 是 E 上一个非退化对称双线性型。

\subsection{平面上的正相似变换。}

回忆（§ 6，第 5 小节），E 的正相似变换是向量空间 E 的一个自同构 u，使得 \(\Phi(u(x), u(y)) = (\det u)\Phi(x, y)\) （对 E 中一切 \(x, y\) ）。

命题 1. --- 设 A(Φ) 是 \(\mathfrak{L}_{\mathbf{A}}(\mathbf{E})\) 的由 E 的正相似变换生成的子代数。

\begin{enumerate}
\def\labelenumi{\alph{enumi})}
\item
  正相似变换就是 A(Φ) 的可逆元素。代数 A(Φ) 是 A 上次数为 2 的交换代数。当 E 不含非零迷向向量时，A(Φ) 是一个域，是 A 的一个二次扩张；否则它是同构于 A 的两个域的直积。
\item
  设 \((e_{1}, e_{2})\) 是 E 的一个正交基；令 \(\alpha_{i} = \Phi(e_{i}, e_{i})\) （i = 1, 2），并令 \(\delta = -\alpha_{2}/\alpha_{1}\) 。A( \(\Phi\) ) 的元素关于这个基的矩阵是形如 \(\begin{pmatrix} a & \delta b \\ b & a \end{pmatrix}\) 的矩阵，其中 \(a \in A\) ，\(b \in A\) 。
\item
  空间 E 是一个自由循环 A(Φ)-模，由任一非迷向向量生成。
\end{enumerate}

事实上，引入 E 上一个辅助的交错双线性型 B ≠ 0；于是 B 非退化。因此存在 E 的一个自同态 w，使得 \(\Phi(x, y) = \mathrm{B}(w(x), y)\) （对 E 中一切 x、y）。对 E 的每个可逆自同态 u，我们有

\[
\begin{array}{r l} \Phi (u (x), u (y)) & = \mathrm{B} (w u (x), u (y)) = \mathrm{B} (u u ^ {- 1} w u (x), u (y)) \\ & = (\det u) \mathrm{B} (u ^ {- 1} w u (x), y); \end{array}
\]

要使 u 是正相似变换，必须且只需

\[
(\det u) \mathrm{B} (u ^ {- 1} w u (x), y) = (\det u) \Phi (x, y) = (\det u) \mathrm{B} (w (x), y)
\]

（对 E 中一切 \(x, y\)）；由于 det \(u \neq 0\) 且 B 非退化，这等价于 \(u^{-1}wu = w\) ，也等价于 \(uw = wu\) 。取 B 为其矩阵 \(S\) 关于 \((e_1, e_2)\) 为 \(\begin{pmatrix} 0 & 1 \\ -1 & 0 \end{pmatrix}\) 的那个交错双线性型；以 \(R\) 与 \(W\) 分别记 \(\Phi\) 与 \(w\) 关于这个基的矩阵，则关系 \(\Phi(x, y) = \mathrm{B}(w(x), y)\) 依 § 1 第 10 小节的公式 (47) 写成 \(R = ^t W.S\) ；把它明写出来即得 \(W = \begin{pmatrix} 0 & \alpha_2 \\ -\alpha_1 & 0 \end{pmatrix}\) 。若以 \(\begin{pmatrix} a & c \\ b & d \end{pmatrix}\) 记 \(u\) 关于 \((e_1, e_2)\) 的矩阵，则关系 \(uw = wu\) 就等价于关系 \(b\alpha_2 = -c\alpha_1\) ，\(a\alpha_2 = d\alpha_2\) ，\(a\alpha_1 = d\alpha_1\) ，也就是 \(a = d\) 与 \(c = \delta b\) ；这证明了正相似变换的矩阵是形如 \(\begin{pmatrix} a & \delta b \\ b & a \end{pmatrix}(a, b\) 在 A 中)的可逆矩阵。

现在，其矩阵形如 \(\begin{pmatrix} a & \delta b \\ b & a \end{pmatrix}\) （a, b 在 A 中）的 E 的自同态构成 \(\mathfrak{L}_{\mathrm{A}}(\mathrm{E})\) 的一个维数为 2 的向量子空间，它由 1 与自同态 w 生成；由于 \(w^{2}\) 是比为 \(-\alpha_{1}\alpha_{2}\) 的位似，这个子空间就是子代数 \(\mathrm{A}(\Phi)\) ，即 \(\mathfrak{L}_{\mathrm{A}}(\mathrm{E})\) 中由正相似变换生成的那个子代数。正相似变换是 \(\mathrm{A}(\Phi)\) 的可逆元素，也就是其矩阵满足 \(a^{2}-\delta b^{2}\neq0\) 的那些元素。代数 \(\mathrm{A}(\Phi)\) 是交换的这一事实是显然的。对它应用第二章 § 7 第 7 小节的结果：若 δ 不是 A 中的平方，即 E 的非零向量都不是迷向的，则 \(\mathrm{A}(\Phi)\) 是一个域；反之，若 δ 是 A 中的平方，即 E 含有非零迷向向量，则 \(\mathrm{A}(\Phi)\) 是同构于 A 的两个域的直和。这就证明了 a) 与 b)。

最后，E 的任一非迷向向量都可以取作

某个正交基的第一个向量 \(e_1\) ，设该正交基为 \((e_1, e_2)\) （ \(\S 6, n^0 1\) ）；因此它的像 \(u(e_1)\) （ \(u\) 为 A( \(\Phi\) ) 的元素）都是形如 \(ae_1 + be_2\) （ \(a, b\) 在 A 中）的向量，也就是 E 中所有的向量，因为所有矩阵 \(\begin{pmatrix} a & \delta b \\ b & a \end{pmatrix}\) 都是 A( \(\Phi\) ) 的元素的矩阵。换言之，E 是一个循环 A( \(\Phi\) )-模，由任一非迷向向量生成。此外我们看到，它是单生成的自由 A( \(\Phi\) -模，因为 \(u(e_1) = ae_1 + be_2 = 0\) 蕴含 \(a = b = 0\) ，从而 \(u = 0\) 。这就证明了 \(c\) )。

注. --- 1) 设 \(\vartheta\) 是矩阵为 \(\begin{pmatrix}0 & \delta \\ 1 & 0\end{pmatrix}\) （关于 \((e_1, e_2)\) ）的相似变换，下面把它记作 v；命题 1 的证明中引入的相似变换 \(w\) 等于 \(-\alpha_1 v\) ；我们有 \(v^2 = \delta\) 。正相似变换 \(u = a + b v\) （ \(a, b\) 在 A 中）的乘子等于它的矩阵 \(\begin{pmatrix}a & \delta b \\ b & a\end{pmatrix}\) 的行列式，即 \(a^2 - \delta b^2 = (a + b v)(a - b v) = u.\bar{u}\) ，其中以 \(\bar{u}\) 记 \(u\) 在代数 A(Φ) 中的共轭元（第二章 § 7，第 7 小节）；换言之，\(u\) 的乘子就是 \(u\) 在代数 A(Φ) 中的范数 N( \(u\) ) （出处同前）。特别地，要使正相似变换 \(u\) 是一个旋转，必须且只需 N( \(u\) ) = 1；要使 \(u\) 是位似，必须且只需 \(u \in A^*\) 。

\begin{enumerate}
\def\labelenumi{\arabic{enumi})}
\setcounter{enumi}{1}
\item
  其平方是位似且形如 \(u = a + bv\) （ \(a, b\) 在 A 中，\(b \neq 0\) ）的正相似变换只有位似以及标量倍 \(bv\) （ \(v\) 的），因为 \((a + bv)^{2} = (a^{2} + \delta b^{2}) + 2abv\) 。后者正是把每个向量变成正交向量的向量空间 E 的那些自同构；事实上，这样的自同构的矩阵必然形如 \(\begin{pmatrix} 0 & c \\ d & 0 \end{pmatrix} (c, d\) 在 A 中)，而把向量 \(\lambda e_{1} + \mu e_{2}\) 变成正交向量的条件这时就写成 \(\lambda \mu(d\alpha_{2} + c\alpha_{1}) = 0\) 。
\item
  容易验证，对 E 中的 \(x\) 、\(y\) 及 \(u \in A(\Phi)\) ，有 \(\Phi(u(x), y) =: \Phi(x, \overline{u}(y))\) 。于是，正相似变换 \(u\) 的伴随自同态是正相似变换 \(\overline{u}\) ，即 \(u\) 在 A( \(\Phi\) ) 中的共轭元。
\item
  由于 E 的每个逆相似变换都是一个正相似变换与关于 \(Ae_{1}\) 的对称（变换）的乘积，逆相似变换关于 \((e_{1}, e_{2})\) 的矩阵是形如 \(\begin{pmatrix} a & -\delta b \\ b & -a \end{pmatrix}\) 的矩阵。
\end{enumerate}

此后我们以 S 表示 E 的相似变换群，以 \(S^{+}\) 表示正相似变换群，以 H 表示非零位似的群，并以 \(O^{+}\) 表示旋转群。回忆我们有 \(H \subset S^{+}\) （§ 6，第 5 小节）。

推论 1. --- 正相似变换的群 S \(^{+}\) 是交换的。对 E 中任意非迷向向量 x、y，存在唯一的正相似变换 u 使 y = u(x)。

第一个断言由代数 A(Φ) 是交换的这一事实得出。由于 E 是由 x（相应地，y）生成的自由单生成 A(Φ)-模，存在 A(Φ) 的唯一元素 u（相应地，u′）使得 \(y = u(x)\) （相应地，\(x = u'(y)\) ）；于是 \(x = u(u'(x))\) ，且 \(uu'\) 是恒等映射；这表明 u 可逆，因此是一个正相似变换。

推论 2. --- 旋转群 O \(^{+}\) 是交换的。对 E 中满足 Φ(x, x) = Φ(y, y) ≠ 0 的任意向量 x、y，存在唯一的旋转 u 使 y = u(x)。

第一个断言由推论 1 得出。它还表明，存在正相似变换 \(u\) 且恰好一个使 \(y = u(x)\) ；由于 \(\Phi(u(x), u(x)) = \Phi(x, x)\) ，\(u\) 的乘子等于 1，因此 \(u\) 是一个旋转。

推论 3. --- 群 S \(^{+}\) /H 是交换的。它作用于 E 的非迷向直线的集合上。无论 E 的非迷向直线 D、D′ 如何，都存在且只存在 S \(^{+}\) /H 的一个元素 φ 使 D' = φ(D)。

这由推论 1 以及 H 使 E 的每条直线整体不变这一事实得出。

命题 2. --- 从 O \(^{+}\) 到 S \(^{+}\) /H 的典范同态的核是 \{1, -1\}。

事实上，1 与 -1 是 \(\mathbf{H} \cap \mathbf{O}^{+}\) 的仅有的元素。

命题 3. --- 同态 \(u \to u/\overline{u} = u^{2}/\mathrm{N}(u)\) （\(\mathrm{S}^{+}\) 到自身的）以 H 为核，并通过取商定义一个从 \(\mathrm{S}^{+}/\mathrm{H}\) 到 \(\mathrm{O}^{+}\) 上的同构。

事实上，关系 \(u/\overline{u} = 1\) 等价于 \(u = \overline{u}\) ，也就是 \(u \in A^* = H\) 。于是 H 是 \(u \to u/\overline{u}\) 的核。由于 \(N(u/\overline{u}) = 1\) ，\(u/\overline{u}\) 是一个旋转（注 1）。剩下要证明每个元素 \(v\) （\(O^+\) 中的）都形如 \(u/\overline{u}\) （ \(u \in S^+\) ）。若 \(1 + v\) 可逆，则可取 \(u = 1 + v\) ，因为关系 \(N(v) = v\bar{v} = 1\) 蕴含 \(1 + v = v(1 + \bar{v})\) 。否则有 \(N(1 + v) = (1 + v)(1 + \bar{v}) = 0\) ，也就是说，令 \(v = a + bw\) （ \(a \in A\) ，\(b \in A\) ，\(w^2 = \delta \in A\) ）后，\(2(1 + a) = 0\) ，从而 \(a = -1\) ；但关系 \(a = -1\) 与 \(N(v) = a^2 - \delta b^2 = 1\) 蕴含 \(b = 0\) ，从而 \(v = -1\) ；由于 \(\bar{w} = -w\) ，在这种情形取 \(u = w\) 即可。

当 A(Φ) 是域时，命题 3 是希尔伯特定理的一个特殊情形（第五章 § 11，第 5 小节，定理 3）。

推论. — 设 \(i: \mathrm{O}^{+} \to \mathrm{S}^{+}/\mathrm{H}\) 与 \(d: \mathrm{S}^{+}/\mathrm{H} \to \mathrm{O}^{+}\) 是命题 2 与 3 中定义的同态，并把交换群 \(\mathrm{O}^{+}\) 与 \(\mathrm{S}^{+}/\mathrm{H}\) 写成加法；则有 \(d(i(\theta)) = 2\theta\) （对 \(\theta \in \mathrm{O}^{+}\) ），且有 \(i(d(\varphi)) = 2\varphi\) （对 \(\varphi \in \mathrm{S}^{+}/\mathrm{H}\) ）。

事实上，若 \(\bar{\nu}\) 是一个旋转，则 \(\bar{\nu} = \nu^{-1}\) ，于是 \(d(i(\nu)) = \nu/\bar{\nu} = \nu^{2}\) 。另一方面，若 \(\varphi \in S^{+}/H\) ，则 \(\varphi\) 是某个正相似变换 \(u\) 的模 H 的类，而 \(d(\varphi) = u/\bar{u} = u^{2}/N(u)\) 模 H 同余于 \(u^{2}\) ；由此得第二个公式。

\subsection{平面三角学。}

在这一小节中，我们选取代数 A(Φ) 的一个生成元 w，使得 \(w^{2} \in A\) 。这样的生成元确定到相差一个位似（第 1 小节，注 2），因此 A 的元素 \(w^{2}\) （我们把它记作 δ）确定到模 A 的非零元素的平方所成的乘法子群 (A*)²。

注. — 当 -1 属于所说的模 (A*)² 的类时，通常选取 w 使 \(w^{2} = -1\) ，这就把它确定到相差一个符号。当这个类含 1 而不含 -1 时，通常选取 w 使 \(w^{2} = 1\) ，这同样把它确定到相差一个符号。

这样，每个元素 \(\nu\) （\(S^{+}\) 的）都可以用一种且仅一种方式写成

\[
\nu = c _ {w} (\nu) + s _ {w} (\nu) \varpi\tag{1}
\]

其中 \(c_w(v)\) 、\(s_w(v)\) 属于 A；元素 \(s_w(v)/c_w(v)\) （属于射影域 \(\tilde{\mathbf{A}}\) ；第二章，2 \(^{e}\) 版，附录 III，第 5 小节）记作 \(t_w(v)\) ；它只依赖于 \(v\) 模 H 的类；于是 \(t_w\) 通过取商定义了一个从 S \(^{+}\) /H 到射影域 \(\tilde{\mathbf{A}}\) 中的映射，约定仍把它记作 \(t_w\) （滥用记号）。我们常以 \(c\) 、\(s\) 与 \(t\) 代替 \(c_w\) 、\(s_w\) 与 \(t_w\) 。我们有 \(c_{-w} = c_w\) ，\(s_{-w} = -s_w\) 与 \(t_{-w} = -t_w\) 。

命题 4. — a) 当 \(w^2 = \delta\) 不是 A 中的平方时（即当 E 不含迷向直线时），从 S⁺/H 到 \(\tilde{\mathbf{A}}\) 的映射 t 是双射。

\begin{enumerate}
\def\labelenumi{\alph{enumi})}
\setcounter{enumi}{1}
\item
  当 δ 是 A 中元素 γ 的平方时，映射 t 是从 S⁺/H 到 \(\tilde{\mathbf{A}}\) 除去元素 1/γ 与 -1/γ 后所成的集合上的双射。
\item
  把 \(\mathrm{S}^{+}/\mathrm{H}\) 写成加法，则对 \(\varphi, \varphi'\) （\(\mathrm{S}^{+}/\mathrm{H}\) 中的），有
\end{enumerate}

\[
t (\varphi + \varphi^ {\prime}) = (t (\varphi) + t (\varphi^ {\prime})) / (1 + \delta t (\varphi) t (\varphi^ {\prime}))\tag{2}
\]

当 \(t(\varphi)\) 与 \(t(\varphi')\) 有限且 \(1 + t(\varphi)t(\varphi')\) \(\neq 0\) 时。

事实上，由于 \(S^{+}/H\) 是把 \(A(\Phi)\) 看作 A 上向量平面时去掉 0 的那些直线所成的集合，t 是单射。另一方面，要使元素 \(a + bw\) （ \(a \in A\) ，\(b \in A\) ，属于 \(A(\Phi)\) ）是一个正相似变换，必须且只需它可逆，即 \(N(a + bw) = a^{2} - \delta b^{2} \neq 0\) ，亦即 \((b/a)^{2} \neq 1/\delta\) ；这就证明了 a) 与 b) 中的满射性断言。最后，相似变换 \(1 + t(\varphi)w\) 与 \(1 + t(\varphi')w\) 的乘积是相似变换 \(1 + \delta t(\varphi)t(\varphi') + (t(\varphi) + t(\varphi'))w\) ，这就证明了 c)。

命题 5. --- 把 O \(^{+}\) 写成加法。对 O \(^{+}\) 的每对元素 θ、θ'，我们有

\begin{enumerate}
\def\labelenumi{(\arabic{enumi})}
\setcounter{enumi}{2}
\tightlist
\item
\end{enumerate}

\[
c (\theta) ^ {2} - \delta s (\theta) ^ {2} = 1\tag{4}
\]

\[
c (\theta + \theta^ {\prime}) = c (\theta) c (\theta^ {\prime}) + \delta s (\theta) s (\theta^ {\prime})\tag{5}
\]

\[
s (\theta + \theta^ {\prime}) = s (\theta) c (\theta^ {\prime}) + c (\theta) s (\theta^ {\prime}).
\]

关系 (3) 确实表达了 \(\mathrm{N}(c(\theta)+s(\theta)\omega)=1\) 。要证明 (4) 与 (5)，只需在 \(\mathrm{A}(\Phi)\) 中计算旋转 \(c(\theta)+s(\theta)\omega\) 与 \(c(\theta')+s(\theta')\omega\) 的乘积。

命题 6. — 设 d 是命题 3 中定义的从 S+/H 到 O+ 上的同构。对 S+/H 的每个使 t = t(φ) 为有限的元素 φ，有

\[
s (d (\varphi)) = 2 t / (1 - \delta t ^ {2}), \quad c (d (\varphi)) = (1 + \delta t ^ {2}) / (1 - \delta t ^ {2}).\tag{6}
\]

事实上，\(\varphi\) 是相似变换 \(1 + tw\) 模 H 的类，于是旋转 \(d(\varphi)\) 就是 \((1 + tw)^2/\mathrm{N}(1 + tw) = (1 + \delta t^2 + 2tw)/(1 - \delta t^2)\) （命题 3），这就证明了 (6)。

推论. --- 对 O⁺ 的每个使 \(t(\theta)\) 为有限的元素 θ，我们有

\[
s (2 \theta) = 2 t (\theta) / (1 - \delta t (\theta) ^ {2}), \quad c (2 \theta) = (1 + \delta t (\theta) ^ {2}) / (1 - \delta t (\theta) ^ {2}). \tag {7}
\]

对每个元素 \(\varphi\) （\(S^{+}/H\) 中的），只要 \(t(\varphi)\) 有限且 \(1 + \delta t(\varphi)^{2} \neq 0\) ，就有

\[
t (2 \varphi) = 2 t (\varphi) / (1 + \delta t (\varphi) ^ {2}).\tag{8}
\]

事实上，这由命题 6 及命题 3 的推论立即得出。

注. --- 公式 (6) 对 \(t = \infty\) 仍然成立，只要把出现在右端的那些有理函数换成它们到射影域的典范开拓，这里射影域指 \(\tilde{\mathbf{A}}\) （第二章，2 \(^{\text{e}}\) 版，附录 III，第 5 小节）；事实上，若 \(t = \infty\) ，则 \(\varphi\) 是 \(w\) 的类，且 \(d(\varphi) = -1\) ，\(s(d(\varphi)) = 0\) ，\(c(d(\varphi)) = -1\) ；而这些正是右端各式在 \(t = \infty\) 时由其典范开拓取到的值。当 \(t(\theta) = \infty\) 时 (7) 同样成立，当 \(t(\varphi) = \infty\) 或 \(1 + \delta t(\varphi)^{2} = 0\) 时 (8) 亦然。类似地，当 \(t(\varphi)\) 与 \(t(\varphi')\) 中只有一个（例如 \(t(\varphi)\) ）为无穷时，公式 (2) 仍成立，只要把它的右端看作单独 \(t(\varphi)\) 的有理函数：事实上，相似变换 \(1 + t(\varphi')w\) 与 \(w\) 的乘积是 \(\delta t(\varphi') + w\) ，而 (2) 的右端的典范开拓所取的值是 \(1/\delta t(\varphi')\) 。最后，当 \(t(\varphi)\) 与 \(t(\varphi')\) 都有限且有 \(1 + \delta t(\varphi)t(\varphi') = 0\) 时，有 \(t(\varphi) + t(\varphi') = 0\) （否则 \(t(\varphi)^{2}\) 将等于 \(1/\delta\) ，这是不可能的（命题 4））；因此可以约定 (2) 的右端的值为 \(\infty\) ，而这个值确实就是左端的值。当 \(t(\varphi)\) 与 \(t(\varphi')\) 都是无穷时，(2) 的右端没有定义。

\subsection{角。}

在本小节及下一小节中，我们假设 A 是一个有序域，从而特征为零。回忆（第六章勘误），若 F 是 A 上一个向量空间，则关系“存在 λ \textgreater{} 0 使 y = λx”是 F - \{0\} 的元素 x、y 之间的一个等价关系，这个关系的每个等价类称为以 0 为原点的开半直线，而开半直线与 \{0\} 的并称为以 0 为原点的闭半直线（或简称半直线）；若 D 是一条直线而 Δ 是含于 D 的一条闭半直线，则 D 是 Δ 与 -Δ 的并，且不含其他闭半直线。我们称一条半直线是迷向的，如果包含它的直线是迷向的。

也回忆（出处同前），给定 A 上有限 n 维的向量空间 F，F 的一个定向就是选取空间 \(\bigwedge^{n}F\) 的两条半直线之一；属于这条半直线的 n-向量称为正的。赋有定向的有限维向量空间称为定向的。

E 的比 \textgreater0 的位似显然构成 H 的一个指数为 2 的子群；我们把它记作 H \(^{+}\) 。典范同态 i : O \(^{+}\) → S \(^{+}\) /H （参见命题 2）是 S \(^{+}\) /H \(^{+}\) 到 S \(^{+}\) /H 上的典范同态与 O \(^{+}\) 到 S \(^{+}\) /H \(^{+}\) 中的典范同态的复合；由于 O \(^{+}\) ∩ H \(^{+}\) = \{1\} ，后一同态是单射。

命题 7. --- 假设 A 是极大有序域且 \(\Phi\) 是正型（§ 7）。则从 O \(^{+}\) 到 S \(^{+}\) /H \(^{+}\) 中的典范同态以及从 O \(^{+}\) /\{1, -1\} 到 S \(^{+}\) /H 中的典范同态都是双射，且 S \(^{+}\) 同构于 O \(^{+}\) × H \(^{+}\) 。

我们已经看到所说的同态是单射，故只需证明第一个是满射。设 \((e_{1}, e_{2})\) 是 E 的一个标准正交基，并设 w 是使 \(w(e_{1}) = e_{2}\) 的正相似变换（第 1 小节命题 1 的推论 1）；于是 \(w^{2} = -1\) （命题 1，b)）。对任一正相似变换 \(u = a + bw\) （ \(a \in \mathbf{A}, b \in \mathbf{A}\) ），我们有 \(\mathrm{N}(u) = a^2 + b^2 > 0\) ，且存在一个且仅一个旋转，它含于 \(\mathbf{A}(\Phi)\) 的与 \(u\) 相同的半直线中，即 \((a^2 + b^2)^{1/2} u\) 。证毕。

推论. --- 给定两条以 0 为原点的半直线 D、D'，存在一个且仅一个旋转 v 使 v(D) = D'。

这些假设确实蕴含 E 不含迷向直线。于是我们的断言由第 1 小节命题 1 的推论 1 得出。

此后我们假设 A 是极大有序域，且型 \(\Phi\) 是正的。在 E 的直线（相应地，以 0 为原点的半直线）的偶 \((\mathrm{D}_1, \mathrm{D}_2)\) 所成的集合中，关系“存在正相似变换（相应地，旋转）u 使 \(u(\mathrm{D}_1) = \mathrm{D}_1'\) 且 \(u(\mathrm{D}_2) = \mathrm{D}_2'\) ”是偶 \((\mathrm{D}_1, \mathrm{D}_2)\) 与 \((\mathrm{D}_1', \mathrm{D}_2')\) 之间的一个等价关系。偶 \((\mathrm{D}_1, \mathrm{D}_2)\) 的等价类按定义称为直线（相应地，半直线）\(\mathrm{D}_1, \mathrm{D}_2\) （依此顺序）的角；我们把它记作 \((\widehat{\mathrm{D}_1, \mathrm{D}_2})\) 。

命题 8. --- 假设 A 是极大有序域，且型 \(\Phi\) 是正的。设 \(D_{1}, D_{2}, D_{1}^{\prime}, D_{2}^{\prime}\) 是 E 中四条以 0 为原点的直线（相应地，半直线）。要使角 \((\widehat{D_{1}, D_{2}})\) 与 \((\widehat{D_{1}^{\prime}, D_{2}^{\prime}})\) 相等，必须且只需角 \((\widehat{D_{1}, D_{1}^{\prime}})\) 与 \((\widehat{D_{2}, D_{2}^{\prime}})\) 相等。

我们来证明所述条件的必要性。设 \(u\) 是使 \(u(\mathrm{D}_{1}) = \mathrm{D}_{1}^{\prime}\) 且 \(u(\mathrm{D}_{2}) = \mathrm{D}_{2}^{\prime}\) 的正相似变换（相应地，旋转）。由命题 1 的推论 3，存在正相似变换（相应地，由命题 7 的推论，存在旋转）\(\nu\) 使 \(\nu(\mathrm{D}_{1}) = \mathrm{D}_{2}\) 。由于群 \(\mathrm{S}^{+}\) （相应地，\(\mathrm{O}^{+}\) ）是交换的，我们有 \(\mathrm{D}_{2}^{\prime} = u(\nu(\mathrm{D}_{1})) = \nu(u(\mathrm{D}_{1})) = \nu(\mathrm{D}_{1}^{\prime})\) ，这表明 \((\widehat{\mathrm{D}_{1}, \mathrm{D}_{1}^{\prime})} = (\widehat{\mathrm{D}_{2}, \mathrm{D}_{2}^{\prime}})\) 。把 \(\mathrm{D}_{2}\) 与 \(\mathrm{D}_{1}^{\prime}\) 互换，充分性即由必要性得出。

由命题 8 得出，对 E 中以 0 为原点的直线（相应地，半直线）的每个角 \((\widehat{\mathrm{D}_1,\mathrm{D}_2})\) ，典范地关联着 \(\mathrm{S}^{+}/\mathrm{H}\) （相应地，\(\mathrm{O}^{+}\) ）中一个完全确定的元素，即正相似变换 \(\nu\) （相应地，旋转 \(\nu\) ）的模 H 的类，这些变换满足 \(u(D_1) = D_2\) （对角的任一代表 ( \(D_1, D_2\) )，角为 ( \(\widehat{D_1, D_2}\) )）。这样我们定义了一个典范双射 \(h\) （相应地，\(h'\) ），它从直线角的集合 \(\mathfrak{A}_0\) （相应地，半直线角的集合 \(\mathfrak{A}\) ）映到 \(S^+/H\) （相应地，\(O^+\) ）上；特别地，对每个 \(\varphi \in \mathfrak{A}\) ，称 \(h(\varphi)\) 为角为 \(\varphi\) 的旋转。对于 \(\mathfrak{A}_0\) 与 \(\mathfrak{A}\) ，我们将通过 \(h^{-1}\) 与 \(h'^{-1}\) 把 \(S^+/H\) 与 \(O^+\) 的交换群结构转移到它们上，并把这样得到的群 \(\mathfrak{A}_0\) 与 \(\mathfrak{A}\) 写成加法。若以 D、\(D', D''\) 表示 E 中以 0 为原点的直线（相应地，半直线），则依定义有

\[
\widehat {(\mathbf {D} , \mathbf {D} ^ {\prime \prime})} = \widehat {(\mathbf {D} , \mathbf {D} ^ {\prime})} + \widehat {(\mathbf {D} ^ {\prime} , \mathbf {D} ^ {\prime \prime})}\tag{9}
\]

（沙勒关系）；

由此推出

\[
\widehat {(\mathrm{D} , \mathrm{D})} = 0, \quad \widehat {(\mathrm{D} , \mathrm{D} ^ {\prime})} = - (\widehat {\mathrm{D} ^ {\prime} , \mathrm{D}}).\tag{10}
\]

注. --- 1) E 中以 0 为原点的直线（相应地，半直线）的集合 L 是交换群 S+/H（相应地，O+）的一个齐性空间，且恒等元是使 L 的所有元素都不变的唯一算子。因此可以对 L 应用第二章 2 \(^{e}\) 版附录 II 第 1 小节的公式；命题 8 是“平行四边形法则”的一个特例，而公式 (9) 与 (10) 是（出处同前）公式 (2) 的特例。

\begin{enumerate}
\def\labelenumi{\arabic{enumi})}
\setcounter{enumi}{1}
\tightlist
\item
  在直线角的群的定义中，可以用典范同构于 S+/H 的群 O+/\{−1, 1\} 来代替群 S+/H（命题 7）。于是从 O+ 到 O+/\{−1, 1\} 上的典范同态对应于从 A 到 A0 上的一个同态，即把两条半直线 Δ、Δ' 的角对应到分别包含 Δ 与 Δ' 的直线 D、D' 的角的那个同态。*当域 A 是实数域时，群 A 是群 A0 的一个二重覆盖。*
\end{enumerate}

由命题 8，凡形如 \((\widehat{\mathbf{D}^{\prime},\mathbf{D}^{\prime\prime}})\) 的直线（相应地，半直线）的角，其中 \(\mathbf{D}'\) 与 \(\mathbf{D}''\) 正交（相应地，凡

形如 \((\widehat{\mathrm{D}}, -\mathrm{D})\) 的角），都相等：这确实由第 1 小节的注 2 得出（相应地，是显然的）。这个直线（相应地，半直线）的角称为直角（相应地，平角）；它是 \(A_{0}\) （相应地，A）中的 2 阶元素。

命题 9. --- 假设域 A 是极大有序的，且 \(\Phi\) 是正型。对每个整数 \(n>1\) ，元素 \(\theta\) （属于直线角的群 \(\mathfrak{A}_{0}\) （相应地，半直线角的群 \(\mathfrak{A}\) ））中满足 \(n\theta=0\) 的个数等于 \(n\) 。

由于 \(\mathfrak{A}_{0}\) 、\(\mathrm{S}^{+}/\mathrm{H}\) 、\(\mathfrak{A}\) 与 \(\mathrm{O}^{+}\) 同构（命题 3），只需对 \(\mathrm{O}^{+}\) 进行证明，即证明恰好有 \(n\) 个旋转 \(\varrho\) 满足 \(\varphi^{n}=1\) 。由于 A 是极大有序域且 \(\mathrm{A}(\Phi)\) 是 A 的次数为 2 的扩张域（命题 1 a)），\(\mathrm{A}(\Phi)\) 是代数闭域（第六章 § 2，第 6 小节，定理 3）。于是 \(n\) 次单位根（在 \(\mathrm{A}(\Phi)\) 中）构成一个 \(n\) 阶循环群（第五章 § 11，第 1 小节，定理 1）。由于 \(\mathrm{N}(u)=u\overline{u}\geqslant0\) （对一切 \(u\in\mathrm{A}(\Phi)\) ），关系 \(u^{n}=1\) 蕴含 \(\mathrm{N}(u)=1\) ，从而 \(u\) 是一个旋转（第 1 小节，注 1）。这就证明了我们的断言。

推论. --- 直角（相应地，平角）是群 \(\mathfrak{A}_{0}\) （相应地，\(\mathfrak{A}\) ）中唯一的 2 阶元素。

最后我们假设平面 E 是定向的。

引理 1. --- 设 u 是 E 的一个正相似变换；所有形如 \(x \wedge u(x)\) 的二向量都属于 \(\bigwedge^{2}\mathbf{E}\) 的同一条闭半直线。

\(u\) 是位似的情形是平凡的。否则，有 \(x \wedge u(x) \neq 0\) （对一切 \(x \neq 0\) ）；设 \(x, y\) 是 \(E\) 的两个向量（ \(x \neq 0, y \neq 0\) ）；存在 \(v \in S^+\) 使 \(y = v(x)\) ，于是 \(y \wedge u(y) = v(x) \wedge uv(x) = v(x) \wedge v(u(x)) = (\det v)(x \wedge u(x))\) ；取 \(E\) 的一个标准正交基，即见 det \(v\) 是正的（命题 1 b)）；由此得我们的断言。

这样，在两个生成元 \(w\) （属于 \(\mathrm{A}(\Phi)\) ，满足 \(w^{2} = -1\) ）之中，存在一个且仅一个，使得二向量 \(x \wedge w(x)\) 对一切 \(x\in\mathbf{E}\) 都是正的。我们选取这个生成元来定义函数 \(c_{w}, s_{w}\) 与 \(t_{w}\) （第 2 小节）。设 \(h\) 与 \(h'\) 是上面定义的从直线角群 \(\mathfrak{A}_{0}\) 到 S+/H 上以及从半直线角群 \(\mathfrak{A}\) 到 O+ 上的典范双射。复合映射 \(t_{w}\circ h\) （从 \(\mathfrak{A}_{0}\) 到射影域 \(\tilde{\mathbf{A}}\) ），\(c_{w}\circ h'\) 与 \(s_{w}\circ h'\) （从 \(\mathfrak{A}\) 到域 A）分别记作 tg、cos 与 sin，称为正切、余弦与正弦函数。映射 \(\varphi\to1/\mathrm{tg}\varphi\) （从 \(\mathfrak{A}_{0}\) 到 \(\tilde{\mathbf{A}}\) ）记作 cotg，称为余切函数。余弦、正弦、正切与余切这些函数统称为三角函数。复合映射 tg∘p 与 cotg∘p（其中 p 表示从 \(\mathfrak{A}\) 到 \(\mathfrak{A}_{0}\) 上的典范同态（上面的注 2））约定仍记作 tg 与 cotg（滥用记号）。

由于这里有 \(\delta = -1\) ，§ 2 的公式 (2)、(8)、(3)、(4)、(5) 与 (7) 给出

\begin{enumerate}
\def\labelenumi{(\arabic{enumi})}
\setcounter{enumi}{10}
\tightlist
\item
\end{enumerate}

\[
\operatorname{tg} \left(\varphi + \varphi^ {\prime}\right) = (\operatorname{tg} \varphi + \operatorname{tg} \varphi^ {\prime}) / (1 - \operatorname{tg} \varphi \operatorname{tg} \varphi^ {\prime})\tag{12}
\]

\[
\operatorname{tg} (2 \varphi) = 2 \operatorname{tg} \varphi / (1 - \operatorname{tg} ^ {2} \varphi)
\]

（对 \(\varphi, \varphi'\) ，\(\mathfrak{A}_0\) 中的）；

\begin{enumerate}
\def\labelenumi{(\arabic{enumi})}
\setcounter{enumi}{12}
\tightlist
\item
\end{enumerate}

\[
\cos^ {2} \theta + \sin^ {2} \theta = 1\tag{14}
\]

\[
\cos (\theta + \theta^ {\prime}) = \cos \theta \cos \theta^ {\prime} - \sin \theta \sin \theta^ {\prime}\tag{15}
\]

\[
\sin (\theta + \theta^ {\prime}) = \sin \theta \cos \theta^ {\prime} + \cos \theta \sin \theta^ {\prime}\tag{16}
\]

\[
\sin (2 \theta) = 2 \operatorname{tg} \theta / (1 + \operatorname{tg} ^ {2} \theta),
\]

\[
\cos (2 \theta) = (1 - \mathrm{tg} ^ {2} \theta) / (1 + \mathrm{tg} ^ {2} \theta)
\]

（对 \(\theta, \theta'\) ，\(\mathfrak{A}\) 中的）。另一方面，依定义或作为前面公式的直接推论，我们有：

\begin{enumerate}
\def\labelenumi{(\arabic{enumi})}
\setcounter{enumi}{16}
\tightlist
\item
\end{enumerate}

\[
\operatorname{tg} \theta = \sin \theta / \cos \theta , \quad \cot \mathrm{g} \theta = \cos \theta / \sin \theta\tag{18}
\]

\[
1 + \mathrm{tg} ^ {2} \theta = 1 / \cos^ {2} \theta , \quad 1 + \cot \mathrm{g} ^ {2} \theta = 1 / \sin^ {2} \theta
\]

（对 θ ∈ 𝒜）。

给定 E 的两个非零向量 \(x, y\) ，这两个向量（依此顺序）之间的角称为并记作 \((\widehat{x}, y)\) ，即它们所属的半直线的角。对 E 的每个向量 \(x\) ，把 A 的一个元素称为 \(x\) 的长度并记作 \(|x|\) ，这个元素就是 \(\Phi(x, x)^{1/2}\) 。

命题 10. --- 假设域 A 是极大有序的，平面 E 是定向的，且型 \(\Phi\) 是正的。对 E 的每对非零向量 \(x\) 、\(y\) ，我们有 \(a\)

\begin{enumerate}
\def\labelenumi{(\arabic{enumi})}
\setcounter{enumi}{18}
\tightlist
\item
\end{enumerate}

\[
\cos \widehat {(x , y)} = \Phi (x, y) / | x |. | y |\tag{20}
\]

\[
\widehat {\sin (x , y)}. e = (x \wedge y) / | x |. | y |,
\]

其中 e 表示使 \(\Phi_{(2)}(e, e) = 1\) 的正二向量。

事实上，由于向量 \(x' = x / |x|\) 与 \(y' = y / |y|\) 的长度都是 1，存在一个且仅一个旋转 \(\nu\) 使 \(\nu(x') = y'\) （第 1 小节，命题 1 的推论 2）。若令 \(\nu = a + bw\) （ \(a, b\) 在 A 中），则依定义有 \(a = \cos(\widehat{x,y})\) 与 \(b = \sin(\widehat{x,y})\) 。关系 \(y' = \nu(x') = ax' + bw(x')\) 给出 \(\Phi(x', y') = a\Phi(x', x') = a\) ，因为 \(x'\) 与 \(w(x')\) 正交（第 1 小节的注 2）；这就证明了 (19)。另一方面，这个关系也给出 \(x' \wedge y' = bx' \wedge w(x') = b.e\) ，这依据开拓 \(\Phi_{(2)}\) （\(\Phi\) 到 \(\bigwedge^{2} E\) 上的）的定义（§ 1，第 9 小节，公式 (37)）以及 \(w\) 的选取；这就证明了 (20)。

注. --- 3) 给定附属於 E 的仿射平面 L 的两条直线（相应地，半直线）D、D'，把 D 与 D' 的角（记作 (D, D')）称为它们的方向在 E 中所成的角（相应地，E 中与 D 及 D' 对应的以 0 为原点的半直线的角）（第二章，2 \(^{e}\) 版，附录 II，第 1 与第 3 小节）。

\begin{enumerate}
\def\labelenumi{\arabic{enumi})}
\setcounter{enumi}{3}
\tightlist
\item
  设 F 是极大有序域 A 上任一维数的一个向量空间，\(\Psi\) 是 F 上一个正的非退化对称双线性型。给定 F 的两个线性无关的向量 \(x\) 、\(y\) ，设 \(\mathbf{F}'\) 是它们张成的向量平面；我们把 \(x\) 与 \(y\) 的角称为把 \(x\) 与 \(y\) 看作平面 \(\mathbf{F}'\) 的元素时的角；记作 \((x,y)\) 。由 (19)，这个角的余弦由下式给出
\end{enumerate}

\[
\cos \widehat {(x , y)} = \Psi (x, y) / | x |. | y |\tag{21}
\]

\((\text{其中 } |x| = \Psi(x, x)^{1/2}\) 也称为向量 \(x)\) 的长度，因而与 \(F'\) 上所选的定向无关；\((x, y)\) 的正弦与正切在改变 \(F'\) 的定向时变号。给定两个成比例的非零向量 \(x, y\) （\(F\) 的），约定 \((x, y) = 0\) 。
